\chapter{The Drinfeld Center and Bulk Algebras}
\label{ch:drinfeld-center}

The Drinfeld center $\cZ(\cC)$ of a monoidal category $\cC$ is the universal braided monoidal category receiving a forgetful functor from~$\cC$.  In the CY programme, it plays two roles: (i)~for $d = 3$, where the topological $E_3$-structure gives a trivial braiding at the topological level ($\pi_1(\mathrm{Conf}_2(\R^3)) = 1$), the quantum group braiding is recovered through $\cZ(\Rep^{E_1}(A))$; (ii)~the Drinfeld center is the categorical incarnation of the chiral derived center $\Zder(A) = C^\bullet_{\mathrm{ch}}(A, A)$ (Hochschild cochains) from Volume~I (Theorem~H), which computes the universal bulk algebra.  The identification $\cZ(\Rep^{E_1}(A)) \simeq \Rep^{E_2}(\Zder(A))$ (Ben-Zvi--Francis--Nadler, Lurie) is the categorical bulk-boundary correspondence.

\section{The Drinfeld center of a monoidal category}
\label{sec:drinfeld-center-def}

\begin{definition}[Drinfeld center]
\label{def:drinfeld-center}
For a monoidal category $\cC$, the \emph{Drinfeld center} $\cZ(\cC)$ is the category of pairs $(X, \beta_{X,-})$ where $X \in \cC$ and $\beta_{X,-} \colon X \otimes (-) \xrightarrow{\sim} (-) \otimes X$ is a half-braiding satisfying the hexagon axiom.  The Drinfeld center is always braided monoidal.
\end{definition}

\section{Drinfeld center as chiral derived center}
\label{sec:drinfeld-chiral-center}

\begin{proposition}
\label{prop:drinfeld-chiral-derived-center}
For an $E_1$-chiral algebra $A$ with monoidal representation category $\cC = \Rep^{E_1}(A)$, the Drinfeld center $\cZ(\cC)$ is equivalent to the representation category of the chiral derived center:
\[
  \cZ(\Rep^{E_1}(A)) \simeq \Rep^{E_2}(Z^{\mathrm{der}}_{\mathrm{ch}}(A)).
\]
This is the bulk-boundary correspondence: the $E_1$ boundary theory determines an $E_2$ bulk theory via the Drinfeld center.
% AP34: The Drinfeld center (= chiral derived center) is the UNIVERSAL BULK.
% It is NOT obtained by bar-cobar inversion (which recovers A itself), nor
% by Verdier duality (which produces B(A!)). The Drinfeld center is
% computed via Hochschild cochains: Z^{der}_{ch}(A) = RHom(Omega(B(A)), A).
% See Vol I, AP25 and AP34.
\end{proposition}

\section{CY enhancement of the Drinfeld center}
\label{sec:cy-drinfeld}

For a CY category $\cC$ of dimension $d$, the Drinfeld center $\cZ(\Rep^{E_1}(\Phi(\cC)))$ inherits a CY structure of dimension $d + 1$ from the Hochschild cochains (the bulk has one extra dimension from the closed-string direction).  This CY enhancement is the categorical analogue of the shifted-symplectic Lagrangian geometry of Volume~I (Theorem~C, ambient complementarity): the boundary theory $A$ determines a Lagrangian in the symplectic bulk $\Zder(A)$.

\section{Quantum group reconstruction from the center}
\label{sec:qg-reconstruction}

The reconstruction problem: given the braided monoidal category $\cZ(\cC)$, recover the quantum group $U_q(\frakg)$ whose representation category it is.  For finite tensor categories, Etingof--Nikshych--Ostrik reconstruction applies.  For the CY programme, the target is to reconstruct the quantum vertex chiral group $G(X)$ from the Drinfeld center of $\Rep^{E_1}(\Phi(\cC))$.  The Tannakian formalism applies when the fiber functor exists; for non-semisimple categories at roots of unity, the modified trace (Geer--Patureau-Mirand--Turaev) replaces the ordinary fiber functor.
